\section{Introduction}

For an irrational real number \(x\), write
\[
 \mu(x)=\sup\left\{\nu>0:
  0<\abs{x-\frac pq}<q^{-\nu}
  \text{ for infinitely many reduced }p/q,\ q\ge2\right\},
\]
allowing the value \(+\infty\). Dirichlet's theorem gives
\(\mu(x)\ge2\). An upper bound concerns every sufficiently large
denominator, not merely a chosen sequence of approximants.
All logarithms of positive real numbers in this paper are natural.

\begin{theorem}[Rational logarithms]\label{thm:rational}
Let \(r\in\Q\), \(r>0\), \(r\ne1\). For every \(\varepsilon>0\)
there is an integer \(Q=Q(r,\varepsilon)\ge2\) such that
\begin{equation}\label{eq:rational-main}
 \abs{\log r-\frac pq}\ge q^{-(2+\varepsilon)}
 \qquad(p\in\Z,\ q\in\N,\ q\ge Q).
\end{equation}
No coprimality condition is imposed. Consequently \(\log r\) is
irrational and
\[
 \mu(\log r)=2.
\]
In particular, \(\mu(\log3)=2\).
\end{theorem}

The broader statement records exactly how the number field enters.
Embeddings, including two complex-conjugate embeddings, are counted
separately.

\begin{theorem}[Compatible embeddings]\label{thm:compatible}
Let \(F\) be a number field of degree \(d=[F:\Q]\), let
\(\gamma,\xi\in F\), \(\gamma\ne0\), and let \(x\in\R\setminus\{0\}\).
Suppose that a nonempty set \(S\) of embeddings \(F\hookrightarrow\C\)
satisfies
\begin{equation}\label{eq:compatibility}
 \exp(\sigma(\gamma)x)=\sigma(\xi)\qquad(\sigma\in S).
\end{equation}
Put \(s=|S|\) and \(\kappa=d/s\).
For every \(\varepsilon>0\) there is an integer \(Q\ge2\) such that
\begin{equation}\label{eq:compatible-main}
 \abs{x-\frac pq}\ge q^{-(2\kappa+\varepsilon)}
 \qquad(p\in\Z,\ q\in\N,\ q\ge Q).
\end{equation}
In particular \(x\) is irrational and
\[
 2\le\mu(x)\le\frac{2d}{s}.
\]
If all embeddings are compatible, then \(\mu(x)=2\).
\end{theorem}

Compatibility already implies \(\xi\ne0\). There is no hypothesis that
\(\xi\) is not a root of unity, that \(x\) is already irrational, or
that an interpolation matrix has full rank. Those facts or alternatives
needed in the proof are established below.

\begin{corollary}[Positive algebraic bases]\label{cor:algebraic}
If \(\alpha>0\) is algebraic over \(\Q\), \(\alpha\ne1\), and
\(d=\deg_{\Q}(\alpha)\), then \(\log\alpha\) is irrational and
\[
 2\le\mu(\log\alpha)\le2d.
\]
For every \(\varepsilon>0\), the associated lower bound with exponent
\(2d+\varepsilon\) holds for all sufficiently large denominators.
\end{corollary}
\begin{proof}
Take \(F=\Q(\alpha)\subset\R\), \(\gamma=1\),
\(\xi=\alpha\), \(x=\log\alpha\), and let \(S\) contain its given real
embedding into \(\C\). Theorem~\ref{thm:rational} is the case
\(F=\Q\), \(s=d=1\).
\end{proof}

\subsection{Prior work and the contribution here}

Quantitative approximation to logarithms has a substantial history.
Zudilin's survey \cite{Zudilin2004} discusses integral and hypergeometric
methods for \(\log2\), \(\log3\), and \(\pi\).
For \(\log3\), Wu and Wang \cite[Corollary~1]{WuWang2014} obtained
\(\mu(\log3)\le5.1163051\), subsequently improved by Bondareva,
Luchin, and Salikhov \cite[Corollary~1]{BLS2018} to
\(\mu(\log3)\le5.116201\). These are finite upper bounds,
not the exact value two asserted in Theorem~\ref{thm:rational}.
Marcovecchio and Viola \cite{MV2012} study approximation to logarithms
by elements of number fields, as well as nonquadraticity measures,
using double integrals and the saddle-point method.
Their number-field approximation problem, normalized by Weil height,
must be distinguished from the rational-approximation exponent used here.
Marcovecchio's multiple-Legendre construction \cite{Marcovecchio2014}
also treats approximation by algebraic numbers of bounded degree.
Bouchelaghem, He, Li, and Wu \cite{BHLW2024} study linear-independence
measures involving two rational logarithms, another quantitative problem
distinct from the individual exponents considered here.
We do not claim that finite irrationality bounds for logarithms were
previously unknown, nor that our degree bound is best in every special
case.
The immediate input is OpenAI's preprint \cite{OAIpi} and its pinned
formal development. Its proof of
\(\mu(\pi)=2\) combines separated approximation denominators, weighted
interpolation, arithmetic clearing of a nonzero determinant, and an
analytic saving caused by repeated Taylor coefficient rows.
This paper reuses that architecture and its proved general ingredients.
It is not a deduction from the numerical statement \(\mu(\pi)=2\).

The new geometric input allows the multiplicative coordinates of the
interpolation points to be distinct nonzero numbers. For \(\log3\)
they are \(3^j\); for \(\log r\), they are \(r^j\).
The local coordinate \(t=Y/y_j-1\) preserves the logarithmic frame.
Once the rigidity argument forces \(Y\) to be constant, such a curve
can meet at most one center. The formal development also supplies the
common-fiber alternative needed when a root of unity is reduced to \(1\).
Both alternatives lead to the actual, explicitly defined interpolation
matrix, not to a matrix of assumed rank.

The arithmetic extension clears denominators of both the algebraic base
\(\xi\) and the slope \(\gamma\). A single selected determinant over
\(F\) is then mapped through every embedding. Compatible embeddings
give analytic decay; the others incur a controlled growth cost.
Taking its norm and dividing by the number of compatible embeddings
produces \(\kappa=d/s\). Omitting the other embeddings would erase this
essential factor and would not prove the stated theorem.

\subsection{Organization and scope}

Section~\ref{sec:preliminaries} fixes the approximation conventions.
Section~\ref{sec:interpolation} proves the required interpolation
statement, making the inherited multiplicity input explicit.
Sections~\ref{sec:arithmetic} and \ref{sec:analysis} establish the
arithmetic and analytic inequalities for the same determinant.
Section~\ref{sec:parameters} chooses all parameters in their required
order and closes the contradiction.
Section~\ref{sec:consequences} gives the specializations, a comparison
with the \(\pi\) proof, and counterexamples to unrestricted real-input
claims. Section~\ref{sec:verification} records formal statements,
source provenance, and reproducible checks.

The results do not assert exponent two for every positive algebraic
base or a uniform finite bound for arbitrary positive real bases.
The eventual thresholds are existential; this paper does not compute
their numerical values.
